% =====================================================================
%  通用问题章模板 —— 每问固定六项子结构（对齐 paper_structure.md §3.2 蓝图）
%  ---------------------------------------------------------------------
%  本文件是“模板的模板”，不直接编译。对应每一问，把它复制为
%  templates/sections/0X_problemN.tex（X = 该问的章节序号+1，N = 问题号），
%  如问题一 → 02_problem1.tex，问题二 → 03_problem2.tex。
%  main_template.tex 的章节树会 \input 这些文件。
%  写完必须回头对照 blueprints/paper_blueprint.md 六项，缺一项即未兑现。
% =====================================================================

% ---------- 递进承接句（必填，每问开头） ----------
% “问题 X 在问题 Y 的基础上，从 …… 提升到了 ……，本问……”（让评委感到递进链）
\section{基于问题N的××模型}   % 主节标题，编号由 ctexset 自动生成

% ---------- 子结构 1：模型 ----------
\subsection{模型建立}
[闭式目标函数（引用规格书）、决策变量、关键假设。判据/核心定义先给直观动机，再给数学形式。]

% ---------- 子结构 2：求解 ----------
\subsection{求解方法}
[实际算法/求解器、参数、是否收敛；多起点数/种子数/网格；如何保证全局性或交叉验证。]

% ---------- 子结构 3：结果 ----------
\subsection{结果与分析}
[关键数字表（booktabs）+ 一句话解读。数字含义，不是数字堆砌。]

% ---------- 子结构 4：策略规律（必填，A196 对比回灌） ----------
\subsubsection{策略规律提炼}
[从最优解回提炼**可迁移的规律**——最优解的几何/物理直觉是什么？对一般情形有什么指导？
范例句式：“最优云团位于导弹进近走廊 x_c≈1.74万 m 处，微小的航向偏角是对准圆柱轮廓而非中心点的精细调整。”
只有数字没有规律的段落不合格。]

% ---------- 子结构 5：亮点 ----------
\subsubsection{本问亮点}
[本问值得强调的一点。]

% ---------- 子结构 6：判据示意（几何/机理类题 ≥1 幅） ----------
% 纯评价/预测类题可只写一句说明（需注明）。
\begin{figure}[H]
\centering
% \includegraphics[width=0.8\textwidth]{figN_scheme.png}
\caption{[判据/几何示意图：先直观后公式]}\label{fig:scheme}
\end{figure}
